\documentclass{article}

% \usepackage[utf8]{inputenc}
% \usepackage[T1]{fontenc}
% \usepackage{lmodern}
% \usepackage[french]{babel}
% \usepackage{userpackages}
% \usepackage{scratch3}
\usepackage{amsmath}
\usepackage{tikz}
\usepackage{enumitem}
\usepackage[a5paper]{geometry}
\usepackage[tikz]{mdframed}
\usepackage{titlesec}
\usepackage{dashrule}

\newgeometry{
    lmargin=50pt,
    hmargin=50pt,
    top=50pt,
    textheight=650pt
}

\titleformat
    {\section}% niveau de titre
    {\LARGE\bfseries}% police numérotation
    {| Partie \Roman{section} |}% texte numérotation
    {8pt}% indentation
    {\normalfont}% police titre

\newcounter{ii}
\long\def\ncar#1#2{%
    \ifnum\arabic{ii}<#1
        #2\stepcounter{ii}\ncar{#1}{#2}%
    \fi
    \setcounter{ii}{0}
}

\def\doted{\hdashrule{.95\linewidth}{.5pt}{2pt}}
\def\dotedlines#1{\ncar{#1}{\vskip10pt\doted}}

\begin{document}\pagestyle{empty}
    \section{Mise en oeuvre d'une configuration de Thalès.}
    \begin{enumerate}[label=\bfseries\arabic*)]
        \item Démarrer le logiciel GeoGebra (par exemple en tapant geogebra
            dans la barre de recherche)
        \item Supprimer les axes et la grille en effectuant un clic droit à
            l'intérieur de la {\sc fenêtre graphique}.
        \item En utilisant les {\sc outils}, tracer dans la {\sc fenêtre
        graphique}, la figure ci-dessous.\par\noindent\vskip10pt
        \begin{mdframed}[
            roundcorner=10pt,
            outerlinewidth=1,
            backgroundcolor=gray!5!
            ]
            \begin{tikzpicture}
                \coordinate[label={[label distance=8pt]175:A}](A)at(0,0);
                \coordinate[label=above left:B](B)at(220:8);
                \coordinate[label=above right:C](C)at(-60:5);

                \path(A)--(B)coordinate[pos=-0.2](A0);
                \path(A)--(B)coordinate[pos=1.3](B0);
                \draw(A0)--(B0);
                \path(A)--(C)coordinate[pos=-0.3](A1);
                \path(A)--(C)coordinate[pos=1.4](C0);
                \draw(A1)--(C0);

                \path(B)--(C)coordinate[pos=-0.3](B1);
                \path(B)--(C)coordinate[pos=1.3](C1);
                \draw(B1)--(C1);

                \path(A)--(B)coordinate[pos=0.3,label=above left:D](D);
                \path(A)--(C)coordinate[pos=0.3,label=above right:E](E);
                \path(D)--(E)coordinate[pos=-1](D0);
                \path(D)--(E)coordinate[pos=2](E0);
                \draw(D0)--(E0);
            \end{tikzpicture}\vskip10pt
            Sur cette figure, les droites (BC) et (DE) sont parallèles. 
            \vskip10pt
        \end{mdframed}
        \item Lever la main pour une vérification du professeur.
    \end{enumerate}
	\newpage
    \section{Triangles semblables}
    Par définition, deux triangles sont semblables s'ils ont les même angles.
    \begin{enumerate}[label=\bfseries\arabic*)]
        \item Que peut-on conjecturer concernant les angles $\widehat{\text{CBA}}$ et
            $\widehat{\text{EDA}}$~?
			\dotedlines2
        \item Que peut-on conjecturer concernant les angles $\widehat{\text{ACB}}$ et
            $\widehat{\text{AED}}$~?
			\dotedlines2
        \item Démontrer les conjectures précédentes.
            \dotedlines6
        \item Que peut-on en déduire concernant les triangles ABC et ADE ?
            \dotedlines6
    \end{enumerate}
    
\end{document}